% !PRINTED 696
\weberpage{696}
Eine Funktion $z$, die nur in je einem Punkte 0 und $\infty$ wird,
also die Ordnung 1 hat, gibt es daher nur in dem Falle $p = 0$.
Jede Funktion in $\Omega$ \gesperrt{ist durch ein solches $z$ rational darstellbar}%
\footnote{Dieser Funktionenkörper führt zu den von den Geometern so
genannten \glqq \gesperrt{Unikursalkurven}\grqq{}.}.

Wenn $r_i > 2$ ist, so sind $\mu_i$ und $z\mu_i$ nach § 198 Differentialquotienten
 erster Gattung. Es gibt also nach (1) Polygone $\mathfrak{W}$, $\mathfrak{W}'$
von der ersten Gattung, die der Bedingung genügen:
\[
\mu_i = \frac{\mathfrak{A}^2\mathfrak{W}}{\mathfrak{Z}},\quad \mu_i z = \frac{\mathfrak{A}^2\mathfrak{W}'}{\mathfrak{Z}},\quad z = \frac{\mathfrak{W}'}{\mathfrak{W}} = \frac{\mathfrak{A}'}{\mathfrak{A}},
\]
und da $\mathfrak{A}$, $\mathfrak{A}'$ teilerfremd sind, so muß $\mathfrak{A}$ in $\mathfrak{W}$, $\mathfrak{A}'$ in $\mathfrak{W}'$ aufgehen.
Es ist also $A$ eine Klasse erster Gattung.

Ist $A$ eine Klasse zweiter Gattung, so sind also alle Exponenten
 $r_1, r_2, \ldots, r_n$ gleich 2 oder kleiner als 2, und es ist im
besonderen $r_{s+1}=2, r_{s+2}=2, \ldots, r_n=2$. Folglich ist
\[
p = (r_2-1)+(r_3-1)+\cdots+(r_n-1) = n-s,
\]
und wir haben nach (2) den Satz:

3. \gesperrt{Die Dimension einer eigentlichen Klasse $A$ von
der zweiten Gattung ist}

\weq{3}{(O,A) = n - p + 1.}

Ist ferner $A$ eine eigentliche Klasse erster Gattung und

\weq{4}{AB = W}

die vollständige Klasse erster Gattung, so nehme man ein Polygon $\mathfrak{B}$
in $B$ und bilde die Differentialquotienten erster Gattung:
\[
w = \frac{\mathfrak{A}^3\mathfrak{B}}{\mathfrak{Z}},\quad zw = \frac{\mathfrak{A}^2\mathfrak{A}'\mathfrak{B}}{\mathfrak{Z}},
\]
und diese Form bleibt erhalten, wenn $\mathfrak{B}$ durch ein äquivalentes
Polygon $\mathfrak{B}'$ ersetzt wird. Die Dimension $(O, B)$ ist also so groß
wie die Dimension der Schar der Differentialquotienten erster
Gattung $w$, die die Eigenschaft haben, daß auch $zw$ noch von
der ersten Gattung ist. Damit dies der Fall sei, dürfen in dem
Ausdruck:
\[
w = y_1\mu_1 + y_2\mu_2 + \cdots + y_n\mu_n
\]
die Grade der Funktionen $y_1, y_2, \ldots, y_n$ nur bis zu der Höhe
$r_1-2, r_2-2, \ldots, r_n-2$, ansteigen, und daraus folgt:
% [sic] printed r_1-2, r_2-2, ..., r_n-2 as the height to which the degrees of y_1, ..., y_n may rise, but that is the bound for w alone to be of the first kind (cf. r_i > 2 above); for zw to be of the first kind as well the degrees can reach only r_i-3, which is what the count (O,B) = (r_{s+1}-2)+...+(r_n-2) on p. 697 presupposes
%CONT>>
